\scanpage{82}\mspage{69}
% archive sheet number: 81

3.11. \unsure{Disons|Dirons} qu'on a une \uline{catégorie à \unsure{cofibration|cofibrations}
stricte} si dans l'axiome $\mathrm{Cof}\,3$ de factorisation
(cf p. 39), la factorisation peut être prise
\deleted{\uline{\unsure{raisonnablement}}}\inserted{fonctoriellement}. Il revient \inserted{$\pm$} au $\unsure{=}$ de dire que
l'axiome de factorisation reste valable,
quand on passe à une catégorie $M(I)$
pour $I$ dans $\mathrm{Cat}$, en y définissant $W(I)$ et
$C(I)$ (à ne pas confondre avec $C_I$) fibre par
fibre. Cela revient au $\mathord{=}$ \deleted{\unsure{tout}} au moins si
$M$ est équiv. à une \unsure{petite} catégorie. Sinon,
il faut supposer fonctorialité pour toute
\uline{petite} sous-catégorie pleine de $\uline{\mathrm{Fl}}\,M$.
Il vaut mieux définir la catégorie à cofibrations
stricte par cette propriété de stabilité.

On voit alors que le prédérivateur
$\mathbb{D}=\mathbb{D}_{\unsure{M|W}}$ \deleted{\unsure{à la}} sur $\mathrm{Cat}$ est \deleted{\illegible{1}} un dérivateur
\uline{relativement à $D=\mathrm{Diag}_0$} défini plus haut
(cat. ordonnées \inserted{$I$} artiniennes, telles que $\mathrm{Esp}(I)$
soit noethérien), ou $\mathord{=}$ rel. à \inserted{$D=$}$\mathrm{Diag}$ (cat
ordonnées $I$ artiniennes sans plus)
lorsque $C$ et $TC$ sont stables par composition
transfinie. J'entends par là que
$u_!$ existe quand $u$ est une flèche de $D$,
